% ===================================================================
%  اللوحة رقم ٢ — جسم الإنسان. --- الفرصة. الألعاب.
%
%  Traduction, faite en 2026, en arabe levantin d'aujourd'hui, sur
%  l'ido de texto/io. Regles de la colonne : voir 10-lawha-01.tex.
%  Une cle porte « pos= », t02-tit-1 : l'imprimeur pose « LA KAPO. »
%  avant l'alinea 1, qui annonce les trois parties du corps ; la
%  traduction garde la cle et sa position.
%
%  LE CORPS EST L'ENDROIT OU LES DEUX REGISTRES SE SEPARENT LE PLUS
%  VISIBLEMENT, ET C'EST LOGIQUE : C'EST CE QU'ON NOMME AVANT
%  L'ECOLE. Le tableau 1 avait montre que les trois arabes du releve
%  se separaient d'abord la ou l'on dit du bien de quelqu'un —
%  منيح contre كويس. Celui-ci montre autre chose : le corps levantin
%  n'est pas un autre vocabulaire, c'est LE MEME, use.
%
%  Deux usures, et elles se voient a chaque alinea. LE ظ EST TOMBE
%  DANS LE ض : le dos (12) est « ضهر » et non ظهر, les os (16) sont
%  « عضام » et non عظام, les ongles (15) sont « ضوافر » et non
%  أظافر. ET LA HAMZA S'EST EFFACEE : la tete est « راس » et non
%  رأس, la main « إيد » et non يد. Rien n'a ete remplace ; deux
%  consonnes ont bouge, et toute la liste change d'aspect.
%
%  A COTE, UNE DIZAINE DE MOTS QUI NE SONT PAS DU TOUT LES MEMES, et
%  ce sont ceux qu'un enfant apprend le premier : « تمّ » la bouche
%  (فم), « منخار » le nez (أنف), « دان » l'oreille (أذن), « زلعوم »
%  la gorge, « بطّة الإجر » le mollet, « إجر » la jambe (رجل),
%  « سنان » les dents, « صوابع » les doigts, « مصارين » les
%  intestins. LA OU LE STANDARD A UN MOT DE LIVRE, LE LEVANTIN A UN
%  MOT DE MERE.
%
%  ET LES JEUX RETOMBENT SUR LE FRANCAIS. La bicyclette (37) est
%  « بسكليت », comme l'abat-jour du tableau 1 etait « أباجورة » : le
%  Mandat a laisse peu de mots, mais il les a laisses sur ce qui
%  roule et sur ce qui s'allume. Les billes (21), elles, sont
%  « جلول » — un mot qui n'est ni standard ni egyptien ni francais,
%  et qu'aucun dictionnaire d'arabe litteraire ne porte. C'est le
%  premier mot de cette colonne qui n'existe que dans une cour de
%  recreation.
%
%  Enfin le cerf-volant (33) est « الطيّارة », l'aeroplane, et c'est
%  le meme mot pour les deux. La langue a nomme l'avion d'apres le
%  jouet, et non l'inverse : le jouet est plus vieux.
% ===================================================================

%%K t02-apar-1 apar
\VUsaut{4.0mm}
\VUtitre{15.0pt}{60}{اللوحة رقم ٢}
\VUsaut{2.0mm}
\VUtitre{11.4pt}{0}{جسم الإنسان. --- الفرصة.}
\VUtitre{11.4pt}{0}{الألعاب.}

%%K t02-tit-0 sub
\VUtitre{13.2pt}{0}{جسم الإنسان.}
\VUsaut{2.4mm}
\VUcentre{10.2pt}{0}{\textit{(درس بسيط بعلوم الطبيعة.)}}
\VUsaut{3.0mm}

%%K t02-tit-1 sub pos=t02-01-1
\VUpk{10.2pt}{40}{الراس.}
\VUsaut{3.0mm}

%%K t02-01-1 p
1. --- أهم أجزاء جسم الإنسان هنّي : \VUgras{الراس} و\VUgras{الجذع}
و\VUgras{الأطراف}.

%%K t02-02-1 p
2. --- الراس فيه جزئين : \VUgras{الجمجمة}، اللي فيها
\VUgras{الدماغ} واللي بيغطّيها \VUgras{الشعر}، و\VUgras{الوجّ}.

%%K t02-03-1 p
3. --- برسمتنا ما منشوف لا الجمجمة ولا الدماغ ؛ بس منشوف كتير واضح
أجزاء الوجّ المهمّة.

%%K t02-04-1 p
4. --- هيّو أولاً \VUgras{الجبين}، اللي بيدلّ على عقل قوي وفكر شغّال
عالدوام. و\VUgras{الصدغين} كتير واسعين. و\VUgras{العين} حيّة ومفتوحة
عالآخر. و\VUgras{حاجب} كتيف و\VUgras{رموش} طويلة عم يحموها.

%%K t02-05-1 p
5. --- وهيّو كمان \VUgras{المنخار} و\VUgras{خروقه}. المنخار بيخدمنا
لنشمّ ولنتنفّس. وعكل جنب من المنخار في \VUgras{الخدود}، وأبعد شوي
\VUgras{الدنين}. وأسفل الوجّ فيه \VUgras{التمّ} و\VUgras{الدقن}.

%%K t02-06-1 p
6. --- ما منشوفن كتير منيح، لأن \VUgras{الشوارب} و\VUgras{السوالف} عم
يخبّوهن عنّا. بس منعرف إنّه التمّ فيه \VUgras{الشفايف} التنتين،
و\VUgras{الفك الأعلى} و\VUgras{الفك الأسفل} مع \VUgras{السنان}،
و\VUgras{اللسان}. بالتمّ منحكي ومناكل. وباللسان مندوق الأكل.

%%K t02-07-1 p
7. --- العين والدان والمنخار والتمّ هنّي، متل الصوابع، أعضاء الحواس
الخمسة : \VUgras{النظر} و\VUgras{السمع} و\VUgras{الشمّ} و\VUgras{الدوق}
و\VUgras{اللمس}.

%%K t02-08-1 p
8. --- فالراس هوّي واحد من أكتر أجزاء جسم الإنسان تشويقاً.

%%K t02-tit-2 sub
\VUsaut{4.4mm}
\VUpk{10.2pt}{40}{الجذع.}
\VUsaut{3.0mm}

%%K t02-09-1 p
9. --- \VUgras{الرقبة} بتوصل الراس بالجذع. وفيها \VUgras{الزلعوم}
و\VUgras{القفا}.

%%K t02-10-1 p
10. --- والجذع كمان فيه أعضاء كتير أساسية. بـ\VUgras{الصدر} في
\VUgras{الرئتين}، اللي منتنفّس فيهن، و\VUgras{القلب} اللي هوّي مركز
الحياة. لمّا القلب بيبطّل يدقّ، الكائن الحي بيموت.

%%K t02-11-1 p
11. --- و\VUgras{الضلوع} بتحمل الصدر.

%%K t02-12-1 p
12. --- وباقي أجزاء الجذع هنّي \VUgras{الجنب} و\VUgras{الضهر}
و\VUgras{الكتاف} و\VUgras{البطن}. منتمدّد عالجنب أو عالضهر لننام،
ومنشيل الحمول على كتافنا.

%%K t02-13-1 p
13. --- وبالبطن في \VUgras{المعدة} و\VUgras{المصارين}. بيخدمونا لنهضم
الأكل.

%%K t02-tit-3 sub
\VUsaut{4.4mm}
\VUpk{10.2pt}{40}{الأطراف.}
\VUsaut{3.0mm}

%%K t02-14-1 p
14. --- عنّا أربع \VUgras{أطراف} : \VUgras{ذراعين} و\VUgras{إجرين}.
بالذراعين مناخد ومنمسك ومنرفع ومنلمّ الأغراض ؛ وبالإجرين منوقف ومنمشي
ومنركض.

%%K t02-15-1 p
15. --- \VUgras{الكتف} بيوصل الذراع بالجسم. و\VUgras{عضلة} كتير قوية،
هيّي العضلة ذات الراسين، بتحرّك الذراع اللي \VUgras{الكوع} بيوصله
بـ\VUgras{الساعد}. و\VUgras{الرسغ} بيعمل مفصل \VUgras{الإيد}، اللي
بتنتهي بـ\VUgras{الصوابع} و\VUgras{الضوافر}. وعالإيد منميّز
\VUgras{عرق} (وشريان) بيجري فيه \VUgras{الدم}.

%%K t02-16-1 p
16. --- وأجزاء الإجر هنّي : \VUgras{الفخد} و\VUgras{الركبة}
و\VUgras{مأخّرها}، و\VUgras{بطّة الإجر} اللي \VUgras{لحمها} مغطّى
بـ\VUgras{الجلد}، و\VUgras{الكاحل} و\VUgras{القدم}. و\VUgras{عضام}
الإجر كتير غليظة وكتير قوية. وبآخر القدم في \VUgras{صوابع الإجر}.
والباقي هوّي \VUgras{الباطن} و\VUgras{الكعب}.

%%K t02-17-1 p
17. --- هيك بيكون الوصف شبه الكامل لجسم الإنسان.

%%K t02-17-2 p
\textit{ملاحظة.} --- \textit{بمساعدة عبارة « عم يوجعني... (راسي وهيك) »
بيقدر الأستاذ يستعمل هالمفردات كلها من جديد من ناحية} \VUgras{وضع
الجسم} \textit{المنيح أو العاطل.}

%%K t02-tit-4 sub
\VUsaut{5.0mm}
\VUpk{10.2pt}{40}{الرياضة.}
\VUsaut{2.4mm}
\VUcentre{10.2pt}{0}{\textit{(بيحكي واحد من التلاميذ.)}}
\VUsaut{3.0mm}

%%K t02-18-1 p
18. --- لنكون مرنين وخفاف وبصحّة منيحة، لازم نمرّن إجرينا وذراعينا.
لهيك منلعب ومنعمل \VUgras{رياضة}.

%%K t02-19-1 p
19. --- خلال الأسبوع، هالتمرينين بيصيروا بساحة الثانوية (شوف اللوحة رقم
١). بس يوم الخميس ويوم الأحد منروح على مرج كبير، وين في محل منركض فيه
ومنتسلّى.

%%K t02-20-1 p
20. --- أساتذتنا وأهالينا بعض المرّات بيجوا معنا لهونيك ؛ بس اللي
بيدير التمارين هوّي \VUgras{أستاذ الرياضة}\textsuperscript{(12)}.

%%K t02-21-1 p
21. --- عالمرج، عشمال المدخل، في \VUgras{أجهزة
الرياضة}\textsuperscript{(1)}، منتسلّق عالـ\VUgras{سلّم}\textsuperscript{(2)}،
ومنعمل الشقلبة عالـ\VUgras{حلقات}\textsuperscript{(3)}
وعالـ\VUgras{ترابيز}\textsuperscript{(4)}، ومنشدّ حالنا بإيدينا لأعلى
\VUgras{سلّم الحبال}\textsuperscript{(5)} أو \VUgras{الحبل
المعقّد}\textsuperscript{(6)}.

%%K t02-22-1 p
22. --- وبنفس الوقت رفقاتنا بينطّوا فوق \VUgras{الحصان
الخشبي}\textsuperscript{(7)} أو \VUgras{المتوازيين}\textsuperscript{(11)}.
وغيرهن \VUgras{بيتلاكموا}\textsuperscript{(8)} أو
\VUgras{بيتبارزوا}\textsuperscript{(9)} بالقناع وبـ\VUgras{الشيش}.
وأخيراً في اللي \VUgras{بيرفعوا الأثقال}\textsuperscript{(10)}.

%%K t02-tit-5 sub
\VUsaut{4.6mm}
\VUpk{10.2pt}{40}{الألعاب.}
\VUsaut{3.0mm}

%%K t02-23-1 p
23. --- وحوالين اللي عم يعملوا رياضة، صبيان وبنات عم يلعبوا
\VUgras{ألعاب} مختلفة. البنات بيتسلّوا بـ\VUgras{الطوق}\textsuperscript{(14)}
تبعهن ؛ وبينطّوا بـ\VUgras{الحبل}\textsuperscript{(13)} أو بيعزموا
إخوتهن وأولاد عمومهن على دورة \VUgras{كروكيه}\textsuperscript{(15)}.
وبينبسطوا لمّا يدفشوا \VUgras{كرة} الخصم بـ\VUgras{مطرقتهن الخشبية}،
بالضبط لمّا يكون مفكّر إنّه رح يمرق من القوس بكل سهولة !

%%K t02-24-1 p
24. --- والصبيان بيفضّلوا الألعاب اللي فيها عضلات أكتر. بيحبّوا يلعبوا
بـ\VUgras{الأوتاد}\textsuperscript{(16)} اللي بيوقّعوها بمهارة
بـ\VUgras{الكرة}\textsuperscript{(17)}. وكمان ما بيستهينوا باللعب
بـ\VUgras{الدوّامة}\textsuperscript{(18)} و\VUgras{البلبوكيه}\textsuperscript{(19)}
ولا حتى بـ\VUgras{لعبة الضفدع}\textsuperscript{(20)}. بس أحلى الألعاب
عندهن هيّي \VUgras{الجلول}\textsuperscript{(21)} و\VUgras{الكرة
عالحيط}\textsuperscript{(22)}، لأن حيط \VUgras{المشتل}\textsuperscript{(26)}
كتير منيح لترجيع الكرة الخفيفة.

%%K t02-25-1 p
25. --- وحنّا الزغير، وهوّي ماشي على \VUgras{ركايزه}\textsuperscript{(24)}،
كتير ماهر وبيعرف يحافظ على توازنه منيح. وجنبه، كم رفيق عم يلعبوا
\VUgras{الغميضة}\textsuperscript{(25)} بهالمشتل ووراء \VUgras{السياج}.
وغيرهن بيفضّلوا \VUgras{لعبة الكفّ}\textsuperscript{(28)} أو
\VUgras{الريشة}\textsuperscript{(29)} اللي بتطير من \VUgras{مضرب} لمضرب.

%%K t02-noto-1 noto
\VUnotes{143.2mm}{%
(*) ألعاب الولاد والصبيان كتير مرّات إلها أسامي وطنية بحتة. نحنا
اجتهدنا لنسمّيها بالإيدو أحسن ما فينا، مرّة مستوحيين من الحركة نفسها
اللي بتميّزها، ومرّة مختارين الاسم الوطني بهالبلد أو هديك، لمّا ما
يكون بعيد كتير. مثلاً : \VUgras{لعبة البرميل} (بالألماني وبالفرنسي)
هيّي \VUgras{لعبة الضفدع} \textsuperscript{(20)} (بالإنكليزي)، اللي
فيها اللاعبين بيحاولوا يرموا أقراصهن المدوّرة من تمّ الضفدع المعدني
اللي واقف فاتح تمّه فوق صندوق (برميل) مخروق.
و\VUgras{القفز عالكتاف}\textsuperscript{(30)} بيفرق عن
\VUgras{القفز عالضهر} بهالشي بس : ليقفزوا فوق الراس، اللاعبين
بيسندوا إيديهن على كتاف رفيقهن، بدل ما يسندوهن — بـ«
\VUgras{القفز عالضهر} » — على ضهره بس. --- أو كمان بعض المشاركين
بيكونوا منحنيين والباقيين رح ينطّوا فوق ضهورهن سانِدين إيديهن
عالكتف ؛ الأوّلين لازم يتحمّلوا الصدمة بلا ما ينثنوا، والتانيين
لازم ينطّوا بلا ما يفقدوا توازنهن (شوف اللوحة نفسها).%
}

%%K t02-26-1 p
26. --- بنص المرج في شجرتين. وعند أصل وحدة منهن، سبعة أو تمن صبيان
نظّموا دورة \VUgras{قفز عالكتاف}\textsuperscript{(30)} (*). مش بكل
البلدان معروفة هاللعبة ؛ بس الكل بيعرف شو هوّي \VUgras{القفز
عالضهر}، اللي الولاد ما عم يلعبوه هالمرة.

%%K t02-tit-6 sub
\VUsaut{5.0mm}
\VUpk{10.2pt}{40}{الألعاب بالهوا الطلق.}
\VUsaut{3.4mm}

%%K t02-27-1 p
27. \VUgras{لعبة الأعمى}\textsuperscript{(31)} كمان كتير مسلّية ؛ البنات
والصبيان بالدور بيحطّوا \VUgras{عصابة العينين} على عيونهن وبيمسكوا
اللي مش شاطرين وبيتركوا حالهن ينمسكوا.

%%K t02-28-1 p
28. --- وبنص المرج، هيّي لعبة كتير فرنسية بس فيها شوي عنف : هيّي
\VUgras{لعبة الدبّ}\textsuperscript{(32)}، وهيّي أكتر ضجّة بكتير من
لعبة \VUgras{الطيّارة}\textsuperscript{(33)} الهادية، وكمان من اللعبة
الإنكليزية \VUgras{التنس}\textsuperscript{(34)}.

%%K t02-29-1 p
29. --- هالرياضة الأخيرة هيّي فرحة التلاميذ الكبار. وأساتذتنا نفسهن
بيمارسوها بكل رغبة. بدّها مهارة كبيرة مع خفّة، وأنا كتير بتعجبني.
ومنترك للبنات ألعاب \VUgras{المرجيحة}\textsuperscript{(35)}.

%%K t02-30-1 p
30. --- ما بتعجبني لعبة إنكليزية تانية، يمكن تصير خطرة.

%%K t02-30-1b p
قصدي \VUgras{كرة القدم}\textsuperscript{(36)} اللي بتفرّح أخوي الكبير.
بفضّل لعبتنا القديمة \VUgras{الأسرى}\textsuperscript{(38)}، اللي هيّي
كمان منيحة للصحّة وأقل عنف بكتير.

%%K t02-tit-7 sub
\VUsaut{4.6mm}
\VUpk{10.2pt}{40}{البسكليت والبالون.}
\VUsaut{3.0mm}

%%K t02-31-1 p
31. --- وكمان بكل رغبة بلفّ حوالين المرج بـ\VUgras{البسكليت}\textsuperscript{(37)}.

%%K t02-32-1 p
32. --- والسنة الجاي رح أتعلّم \VUgras{ركوب الخيل}\textsuperscript{(39)}.
شو ما أحلى حصان عم يخطي أو يهرول برشاقة، وبيقفز فوق الحواجز وهوّي راكض !

%%K t02-33-1 p
33. --- وبالصيف، منتعلّم \VUgras{السباحة}\textsuperscript{(40)}. منغطس
ومنستحمّي بلذّة بميّ النهر الصافية والباردة. وبعض المرّات، بالآخر،
منطيّر \VUgras{بالون}\textsuperscript{(41)} من ورق بيطلع بعظمة بالهوا،
أول ما ينعبّى هوا سخن.

%%K t02-34-1 p
34. --- في كمان ألعاب تانية كتير، بس ما بخلص لو بدّي عدّها، ومن هالملخّص
بيبيّن إنّه بأيام الخميس ما منزهق أبداً على مرجنا الحلو.

%%K t02-34-2 p
ملاحظة. --- \textit{بيقدر الأستاذ يكمّل هالفصل بسهولة، وبيوصف بتفصيل
أكتر كل لعبة من الألعاب الأهم.}
